\section{Foundations: Geometry, Time, the Law, the Dominant Phase}
\label{sec:foundations}

\subsection{The geometry is unique (P5)}

A causal set fixes its continuum geometry through the Hauptvermutung: order plus number
recovers a Lorentzian manifold essentially uniquely. Eight independent sprinklings of
the same region recover the dimension as $3.02 \pm 0.05$, and the proper-time distance
across the region has coefficient of variation $0.027$. \textbf{Where we landed:} the
smooth geometry a discrete order carries is sharp, not just in dimension but in
distance. The conjecture itself remains unproven, here and in the literature.

\subsection{Spacetime is the dominant phase (P2, P6, P12)}

Most causal sets are not manifold-like, the layered Kleitman-Rothschild orders dominate
the uniform measure by sheer number. A theory in which spacetime is real must show its
dynamics defeats them. The sharp Benincasa-Dowker action fails (it drives the ensemble
to layered orders). The \emph{smeared} action fixes it. The decisive tool was a correct
uniform-measure sampler (single-pair toggle, accept only if still transitive),
validated exactly against enumeration, which reaches the large-$N$ entropic regime where
the layered orders take over by $N=64$. There the smeared action makes manifold-like
spacetime a stable phase coexisting with the layered one, a first-order signature
holding across $N=48$ to $160$, and the 2D specialisation (P6) is genuinely
two-dimensional (Myrheim-Meyer dimension $2.0$ to $2.1$). Finally a Wang-Landau
estimate of the density of states (the 1/t schedule) measures the free-energy crossing
(P12): $\beta^\star \approx 0.14$, roughly $N$-independent. \textbf{Where we landed:}
the dynamics does not merely permit smooth spacetime, it makes it the \emph{dominant}
phase of the sum over histories above a finite coupling, with the layered phase
metastable. The remaining nuance is the free-energy crossing at large $N$, which needs
a height-changing cluster move (the single-pair move cannot extend a chain).

\subsection{The law is a trilemma, resolved by the emergent Hamiltonian (P1)}

Is the dynamics governed by a Hamiltonian that is both bounded below and local? For a
reversible rule's own logarithm $H = i \log U$ the answer is a trilemma: local,
bounded-below, and information-propagating cannot all three hold, any two can. The
principal Hamiltonian of a nontrivial reversible rule is nonlocal (range about $0.8N$),
and the explicit local branch exists only for non-propagating commuting-gate rules.
\textbf{Where we landed:} the resolution is to stop using the log of a microscopic step.
The emergent-mesh Hamiltonian, the graph Laplacian or the Dirac operator on the grown
mesh, is local (range $1$, size-independent), bounded below (spectrum from $0$), and
propagating (a finite-speed lightcone) all at once. The rule builds the geometry, and
time is what the local operator does on it. The rule and emergent time are distinct
objects.
